% SHARD-ID: AQM-50-F3-REGULAR-FUNCTIONS-LINE-BY-LINE
% SHARD-TITLE: Regular Functions on \(\mathbb F_3[x]\) Line by Line
% SHARD-SUMMARY: Works through regular functions on \(\mathbf A^1_{\mathbb F_3}\) as residue profiles and finite support selectors.
% SHARD-SUMMARY: Computes explicit reduced smeared Weyl labels for rational and quadratic supports.
% SHARD-SUMMARY: Shows exactly where repeated factors force the Artin-ring thickened layer.
% SHARD-KEYWORDS: F3, regular function, residue profile, smearing, Chinese remainder, Artin ring

\section{Regular Functions on \texorpdfstring{\(\mathbb F_3[x]\)}{F3[x]} Line by Line}
\label{sec:f3-regular-functions-line-by-line}

\subsection{Status and Source Anchors}
This is the explicit \(k=\mathbb F_3\) companion to AQM-49.  It uses AQM-29 for
the point set of \(\Spec\mathbb F_3[x]\), AQM-28 for finite-support Weyl
operators, AQM-34 for polynomial quotients, and AQM-36 for Artin thickenings.

\subsection{Lamport Dependency Skeleton}
\[
\begin{array}{rcl}
D1&:&k=\mathbb F_3,\quad R=k[x],\quad X=\Spec R.\\
D2&:&x_r=(x-r)\quad(r=0,1,2),\quad
      q_0=x^2+1.\\
L1&:&x\text{ as a function has infinite closed support}.\\
L2&:&x^{-1}\in\Gamma(D(x),\mathcal O_X)\text{ has infinite support in }D(x).\\
D3&:&h=x(x-1),\quad Z_{\rm cl}(h)=\{x_0,x_1\}.\\
T1&:&W_h^{\rm red}(F,G)\text{ is a two-qutrit operator determined by
      values at }0\text{ and }1.\\
T2&:&\text{Every two-qutrit label on }Z_{\rm cl}(h)\text{ is represented by
      polynomial pairs modulo }h.\\
D4&:&q_0=x^2+1,\quad K_{q_0}=k[\alpha]/(\alpha^2+1).\\
T3&:&Z_{\rm cl}(q_0)\text{ gives one }9\text{-level qudit}.\\
D5&:&m=x^2q_0^3.\\
T4&:&m\text{ has reduced Hilbert dimension }27\text{ and thickened dimension }
      3^8.\\
\end{array}
\]

\subsection{Step 1: The Space and Its First Closed Points}
Let
\[
  R=\mathbb F_3[x],\qquad X=\Spec R.
\]
The closed points are \(x_\pi=(\pi)\) for monic irreducible \(\pi\).  The
rational closed points are
\[
  x_0=(x),\qquad x_1=(x-1),\qquad x_2=(x-2).
\]
Their residue fields are all \(\mathbb F_3\).  The three monic irreducible
quadratics are
\[
  q_0=x^2+1,\qquad q_1=x^2+x+2,\qquad q_2=x^2+2x+2.
\]
For \(q_0\), write
\[
  K_{q_0}=\mathbb F_3[x]/(q_0)=\mathbb F_3[\alpha],
  \qquad \alpha^2=-1=2.
\]

\subsection{Step 2: A Global Regular Function Is Usually Infinite Support}
Take the regular function \(F=x\in R\).  At a closed point \(x_\pi\),
\[
  F(x_\pi)=x\bmod\pi\in K_\pi.
\]
At rational points,
\[
  F(x_0)=0,\qquad F(x_1)=1,\qquad F(x_2)=2.
\]
At \(x_{q_0}\),
\[
  F(x_{q_0})=\alpha\in K_{q_0}.
\]
\paragraph{Lemma \(L1\).}
The profile of \(x\) is nonzero at every closed point except \(x_0\).

\emph{Proof.}
\(x\bmod\pi=0\) iff \(x\in(\pi)\), equivalently \(\pi\mid x\).  Since \(x\) is
monic irreducible, this happens only for \(\pi=x\).  Thus all closed points
except \(x_0\) lie in the support.

Therefore \(x\) is not a quasi-local Weyl label.  It is a classical residue
profile with cofinite support.

\subsection{Step 3: A Local Regular Function Still Has Infinite Support}
Let
\[
  U=D(x)=X\setminus\{x_0\}.
\]
The function \(x^{-1}\) is regular on \(U\).  Its residue value at \(x_\pi\in U\)
is
\[
  x^{-1}(x_\pi)=(x\bmod\pi)^{-1}\in K_\pi.
\]
At rational points in \(U\),
\[
  x^{-1}(x_1)=1,\qquad x^{-1}(x_2)=2,
\]
because \(2^{-1}=2\) in \(\mathbb F_3\).  At \(q_0\),
\[
  x^{-1}(x_{q_0})=\alpha^{-1}=2\alpha,
\]
since \(\alpha(2\alpha)=2\alpha^2=4=1\).
\paragraph{Lemma \(L2\).}
The profile of \(x^{-1}\) is nonzero at every closed point of \(D(x)\).

\emph{Proof.}
For \(x_\pi\in D(x)\), \(\pi\nmid x\).  Hence \(x\bmod\pi\ne0\) in the field
\(K_\pi\), so it has an inverse.  An inverse is never zero.

Thus local regularity on \(D(x)\) does not make \(x^{-1}\) a finite-support
Weyl label.

\subsection{Step 4: A Polynomial as a Finite Support Selector}
Now take
\[
  h=x(x-1).
\]
The closed zero set is
\[
  Z_{\rm cl}(h)=\{x_0,x_1\}.
\]
The reduced support Hilbert space is
\[
  \mathcal H_h^{\rm red}
  =
  \ell^2(\mathbb F_3)\otimes\ell^2(\mathbb F_3),
\]
so it has dimension \(9\).  The reduced label space is
\[
  E_h^{\rm red}=\mathbb F_3^2\oplus\mathbb F_3^2.
\]

\subsection{Step 5: A Concrete Smeared Two-Qutrit Operator}
Choose
\[
  F=1,\qquad G=x+1.
\]
At \(x_0\),
\[
  (F,G)=(1,1).
\]
At \(x_1\),
\[
  (F,G)=(1,2).
\]
Therefore
\[
  e_h^{\rm red}(F,G)=((1,1),(1,2)).
\]
If \(T_r(a)\) and \(R_r(b)\) are the qutrit shift and phase at \(x_r\), then
\[
  W_h^{\rm red}(F,G)
  =
  T_0(1)R_0(1)\otimes T_1(1)R_1(2).
\]
This is a genuine quasi-local observable because the support has two closed
points.
\subsection{Step 6: No Compatibility Between the Two Qutrits}
Prescribe the label
\[
  q_0=2,\quad q_1=1,\qquad p_0=1,\quad p_1=0.
\]
We seek \(F,G\in\mathbb F_3[x]\) with
\[
  F(0)=2,\quad F(1)=1,\qquad G(0)=1,\quad G(1)=0.
\]
Take
\[
  F=2+2x,\qquad G=1+2x.
\]
Then
\[
  F(0)=2,\quad F(1)=2+2=1,
\]
and
\[
  G(0)=1,\quad G(1)=1+2=0.
\]
\paragraph{Theorem \(T2\).}
Every label in \(\mathbb F_3^2\oplus\mathbb F_3^2\) is obtained by polynomial
pairs modulo \(x(x-1)\).

\emph{Proof.}
The ideals \((x)\) and \((x-1)\) are coprime, since
\[
  1=x-(x-1).
\]
The Chinese remainder theorem gives
\[
  \mathbb F_3[x]/(x(x-1))
  \cong
  \mathbb F_3[x]/(x)\times\mathbb F_3[x]/(x-1)
  \cong \mathbb F_3\times\mathbb F_3.
\]
Apply this once to the position labels and once to the momentum labels.

\subsection{Step 7: A Quadratic Support Is One Nine-Level Qudit}
Take
\[
  h=q_0=x^2+1.
\]
Then
\[
  Z_{\rm cl}(h)=\{x_{q_0}\}.
\]
The residue field is
\[
  K_{q_0}=\mathbb F_3[\alpha]/(\alpha^2+1),
\]
which has \(9\) elements.  Thus the reduced Hilbert space is
\[
  \mathcal H_{q_0}^{\rm red}=\ell^2(K_{q_0}).
\]
Choose
\[
  F=x+1,\qquad G=2x.
\]
Then the reduced label is
\[
  (F\bmod q_0,\ G\bmod q_0)=(\alpha+1,\ 2\alpha)\in K_{q_0}^2.
\]
The smeared operator is the one-site \(9\)-level Weyl operator
\[
  W_{q_0}^{\rm red}(F,G)=W_{q_0}(\alpha+1,2\alpha).
\]
\subsection{Step 8: Repeated Factors and Artin Data}
Take
\[
  m=x^2q_0^3.
\]
Its reduced closed support is
\[
  Z_{\rm cl}(m)=\{x_0,x_{q_0}\}.
\]
Therefore the reduced Hilbert dimension is
\[
  \dim\mathcal H_m^{\rm red}
  =
  |\mathbb F_3|\cdot|K_{q_0}|
  =
  3\cdot9=27.
\]
The exponents \(2\) and \(3\) have disappeared.

In the thickened layer, the local Artin rings are
\[
  A_0=\mathbb F_3[x]/(x^2),\qquad
  A_q=\mathbb F_3[x]/(q_0^3).
\]
Their sizes are
\[
  |A_0|=3^2,\qquad |A_q|=3^{3\deg q_0}=3^6.
\]
Thus
\[
  \dim\mathcal H_m^{\rm thick}
  =
  |A_0|\cdot|A_q|
  =
  3^8=6561.
\]
\paragraph{Theorem \(T4\).}
For \(m=x^2q_0^3\), the reduced theory sees two closed-point qudits, while the
thickened theory sees jet degrees of freedom of total Hilbert dimension
\(3^{\deg m}=3^8\).

\emph{Proof.}
The reduced theory uses only the distinct irreducible factors \(x\) and
\(q_0\), so its dimension is \(3^1\cdot3^2=27\).  The thickened theory uses the
Artin rings \(R/(x^2)\) and \(R/(q_0^3)\), whose sizes are \(3^2\) and
\(3^6\).  Tensoring the two local Hilbert spaces gives dimension \(3^8\).
This equals \(3^{\deg m}\), since \(\deg m=2+3\cdot2=8\).

\subsection{Conclusion}
For \(\mathbf A^1_{\mathbb F_3}\), a regular function such as \(x\) or
\(x^{-1}\) on \(D(x)\) is a residue profile with infinite closed support.  A
nonzero polynomial \(h\) becomes finite only when it is used as an equation:
it selects \(Z_{\rm cl}(h)\).  A Weyl observable then needs labels on that
finite support.  Repeated factors of \(h\) are invisible in the reduced
support net and become visible exactly in the Artin-ring thickened layer.
